\documentclass[12pt, a4paper]{llncs}

\usepackage[backend=biber, style=apa, sorting=ynt]{biblatex}
\usepackage{graphicx}
\usepackage{float}
\usepackage{booktabs}

\title{Energy Efficiency in Railway Scheduling}
\author{Paul Tennigkeit \and Georgi Abshilava}
\institute{University of Potsdam, Germany \\
\email{\{tennigkeit, abshilava\}@uni-potsdam.de}}

\begin{document}
\maketitle

\section{Introduction}
\section{Encoding Energy Efficiency in ASP}
\subsection{What is Flatland?}
Flatland is a framework for simulating railway networks. A Flatland environment consists of a two dimensional plane made up of different types of cells, some of which containing different types of interconnected tracks, and agents moving around on them from a starting point to an end point. Agents have an earliest departure and arrival time, and a pre-defined speed, represented by an inverse value. To limit the time of the simulation, a global time frame is set.

\subsection{Baseline Encoding}
The baseline ASP encoding (train\_encoding.lp) handles all essential train and environment behavior, including defining legal transitions and movements, possible directions, and actions.

\paragraph{Directions}
There are four possible directions agents can go and point towards: North, south, east and west. Each is defined as a fact:
\begin{verbatim}
dir(n). dir(s). dir(e). dir(w).
\end{verbatim}

\paragraph{Movements}
Each possible movement is defined as a fact:
\begin{verbatim}
moving(move_forward).
moving(move_left).
moving(move_right).
\end{verbatim}
\texttt{delta()} defines what it means to "move" in each direction:
\begin{verbatim}
% The origin is located in the top left in Flatland, and 
% the coordinates increment downwards

delta(A,s,1,0)  :- moving(A).
delta(A,n,-1,0) :- moving(A).
delta(A,e,0,1)  :- moving(A).
delta(A,w,0,-1) :- moving(A).
\end{verbatim}

\paragraph{Transitions}
In order to represent different types of tracks, \texttt{transitions()} predicates are defined for each:
\begin{verbatim}
% ID  : track ID
% Dir1: pointing direction of agent
% Act : associated action with movement towards exit
% Dir2: exit direction
transition(ID, Dir1, Act, Dir2).
\end{verbatim}
Due to the way this predicate is built, each track type is defined multiple times depending on its number of exit directions (e.g. straight tracks have two possible exits, north and south, and are, therefore, defined twice for each possible exit).

\paragraph{Actions}
The \texttt{action()} predicate is used as an interface between the Flatland visualization engine and the Clingo logic:
\begin{verbatim}
% train(ID): A train with an ID
% A        : An action
% T        : The time step action A is taken
action(train(ID),A,T) :- chosen(ID,T,A,_,_,_,_,_).
action(train(ID),move_forward,Dep) :- start(ID,_,Dep,_),
                                      Dep > 0.
                                      
action(train(ID),wait,T) :- start(ID,_,Dep,_),
                            Dep > 0,
                            T = 0..Dep-1.
\end{verbatim}

\paragraph{Collision Constraints}
Like in real life, collisions between trains are to be avoided. To achieve this, the following constraints are employed:
\begin{verbatim}
% Trains can not occupy the same cell at the same time
:- stay(ID1,X,Y,T),
   stay(ID2,X,Y,T),
   ID1 < ID2.

% Trains can not swap their position
:- stay(ID1,X1,Y1,T),
   stay(ID2,X2,Y2,T),
   stay(ID1,X2,Y2,T+1),
   stay(ID2,X1,Y1,T+1),
   ID1 < ID2.
\end{verbatim}

\paragraph{Reachability Constraints}
To simulate movement across the grid over time, a train's initial departure parameters are first converted into a \texttt{position()} state:
\begin{verbatim}
position(ID,(X,Y),Dep+1,Dir) :- start(ID,(X,Y),Dep,Dir).
\end{verbatim}

From any current position, valid future states are derived using the \texttt{reachable()} predicate by checking cell track transitions, movement coordinate deltas, and the train's speed:
\begin{verbatim}
reachable(ID,T,A,X2,Y2,T2,Dir2) :- A != wait,
                                   position(ID,(X,Y),T,Dir1),
                                   cell((X,Y),Cell),
                                   transition(Cell,Dir1,A,Dir2),
                                   delta(A,Dir2,DX,DY),
                                   X2 = X + DX,
                                   Y2 = Y + DY,
                                   speed(ID,S),
                                   T2 = T + S,
                                   time(T2).
                                   
reachable(ID,T,wait,X,Y,T+1,Dir) :- position(ID,(X,Y),T,Dir),
                                    time(T+1).
\end{verbatim}

At each discrete time step, a choice rule ensures exactly one valid transition is taken until the destination is reached:
\begin{verbatim}
1 { position(ID,(X2,Y2),T2,Dir2) 
    : reachable(ID,T,A,X2,Y2,T2,Dir2) } 1 
    :- position(ID,(_,_),T,_), not reached(ID,T).
\end{verbatim}

In order to properly enforce the collision constraints, the time interval a train spends inside a specific cell must be tracked:
\begin{verbatim}
chosen(ID,T,A,X,Y,X2,Y2,T2) :- position(ID,(X,Y),T,_),
                               position(ID,(X2,Y2),T2,_),
                               reachable(ID,T,A,X2,Y2,T2,_).
\end{verbatim}

Because trains can have an inverse speed greater than 1 (meaning they traverse a cell over multiple time steps), the \texttt{stay()} predicate marks every time step a cell is occupied between entering and leaving:
\begin{verbatim}
% for inverse speed !=1: define interval to stay in a cell
stay(ID,X,Y,T1) :- chosen(ID,T,_,X,Y,_,_,T2),
                   time(T1), T <= T1, T1 < T2.
\end{verbatim}

\paragraph{Target Arrival}
A train successfully completes its journey when its current coordinates match its designated target destination:
\begin{verbatim}
reached(ID,T) :- position(ID,(X,Y),T,_),
                 end(ID,(X,Y),_). 
\end{verbatim}

To ensure valid and bounded schedules, a look-ahead integrity constraint terminates search branches where a train can physically no longer reach its end coordinates within its individual allowable arrival deadline:
\begin{verbatim}
% stop searching if train is not in time
% to reach its target anymore
:- position(ID,(X,Y),T,_),
   end(ID,(EX,EY),Arr),
   speed(ID,InvSpd),
   T + (|X-EX| + |Y-EY|)*InvSpd > Arr.
\end{verbatim}
By evaluating the remaining Manhattan distance multiplied by the train's inverse speed directly against its scheduled deadline, this single rule simultaneously prevents late arrivals and prunes unviable paths early in the search space.

\paragraph{Schedule Optimization}
Once valid routing pathways are established, the solver optimizes the schedule by minimizing overall transit delays. To ensure that trains are evaluated solely on their actual transit time rather than time spent idling at their destination after completion, the encoding isolates the precise time step $T$ at which a train first reaches its end coordinates:
\begin{verbatim}
% the arrival times should be minimized
reached_earlier(ID,T) :- reached(ID,T),
                         reached(ID,T0),
                         T0 < T.
arrival(ID,T) :- reached(ID,T),
                 not reached_earlier(ID,T).
\end{verbatim}

The \texttt{\#minimize} statement aggregates the arrival times of all agents and assigns this objective a priority level of \texttt{@2}. This high priority ensures that in extended multi-objective encodings minimizing passenger schedule delay always takes precedence over secondary optimization metrics.
\begin{verbatim}
% minimize summed arrival times
#minimize { T@2, ID : arrival(ID,T) }.
\end{verbatim}

\subsection{Energy Encoding Mechanics}
The energy efficiency encoding (train\_encoding\_energy.lp) adds new rules to represent the concept of energy and how different actions of the agents affect their usage and consumption of energy.

\paragraph{Cost Factors}
In real life, trains use a variety of motors that rely on different sources of energy, such as electric, hybrid, diesel or hydrogen power, which fundamentally affects how efficiently energy is consumed. The \texttt{pow\_cost\_fac()} predicate represents this baseline efficiency in the encoding by assigning an integer multiplier to each train based on its power type.

To then determine the true energy footprint of an agent, its mass must also be accounted for. This is handled by the \texttt{cost\_fac()} rule, by multiplying a trains assigned cost factor, the \texttt{PCF}, by its defined weight \texttt{W}, resulting in a final, individualized cost factor \texttt{CF} for each train in the simulation:
\begin{verbatim}
% energy cost factor of trains based on power type + weight
pow_cost_fac(ID,1) :- pow_type(ID, electro).
pow_cost_fac(ID,2) :- pow_type(ID, hybrid).
pow_cost_fac(ID,3) :- pow_type(ID, diesel).
pow_cost_fac(ID,4) :- pow_type(ID, hydrogen).

cost_fac(ID,CF) :- pow_cost_fac(ID,PCF),
                   weight(ID,W),
                   CF = PCF * W.
\end{verbatim}

\paragraph{Energy Costs of Actions}
Depending how fast the train moves or accelerates, it consumes different amounts of energy. These are reflected by the different \texttt{energy()} predicates:

\begin{verbatim}
% This rule calculates the energy cost of waiting,
% which is CF * 24 / S. The penalty is only applied 
% once, when the train first hits the brakes.
% This behavior is reflected by the rule
% not prev_wait(ID,T).
energy(ID,T+1,E) :- action(train(ID),wait,T),
                    speed(ID,S),
                    not prev_wait(ID,T),
                    cost_fac(ID,CF),
                    E = CF * 24 / S.

% This rule calculates the energy cost of moving
% straight. This action's cost is only the base CF.
energy(ID,T2,E) :- action(train(ID),A,T),
                   A != wait,
                   position(ID,_,T,Dir),
                   speed(ID,S),
                   T2 = T + S,
                   position(ID,_,T2,Dir),
                   cost_fac(ID,CF),
                   E = CF .

% This rule calculates the energy cost of turning.
energy(ID,T2,E) :- position(ID,_,T,Dir1),
                   speed(ID,S),
                   T2 = T + S,
                   position(ID,_,T2,Dir2),
                   Dir1 != Dir2,
                   cost_fac(ID,CF),
                   E = CF * 12 / S.
                   
% The prev_wait() rule is required in order to make sure,
% the agent is not constantly penalized for waiting in
% every time step.
prev_wait(ID,T) :- action(train(ID),_,T), 
                   time(T-1),
                   action(train(ID),wait,T-1).
\end{verbatim}

\paragraph{Optimizing the Energy Encoding}
The energy encoding adds the rule \texttt{\#minimize \{ E@1,ID,T : energy(ID,T,E) \}.} to also optimize fuel efficiency. As it is assigned the weight \texttt{@1}, the solver optimizes for arrival time first, which is assigned \texttt{@2}. Fuel efficiency is then optimized within the bounds of the best possible arrival times. 

\subsection{Acceleration Encoding Mechanics}
This is an extension to the energy encoding, that also models the trains' acceleration.

\paragraph{Maximum and Initial Speeds}
The environment created for the experiments also defined the trains' speeds in form of \texttt{speed()} predicates, which is actually interpreted as an inverse speed (e.g. a \texttt{speed} of 4 actually means, that a train would be 25\% slower). To model acceleration, \texttt{speed()} is reinterpreted as \texttt{max\_speed()}.

\begin{verbatim}
% interpret given train speeds as maximum speeds
max_speed(ID,S) :- speed(ID,S).
\end{verbatim}

As a train starts moving, it initially has a speed of 0.25.

\begin{verbatim}
% trains start with a speed of 1/4 
temp_speed(ID, Dep+1, 4) :- start(ID,(_,_),Dep,_).
\end{verbatim}

\paragraph{Dynamic Speed Change}
The \texttt{temp\_speed()} choice rules allow gradual speed changes while strictly respecting a train's physical limits. During each position change, a train can alter its inverse speed by a maximum of one step at a time.

\begin{verbatim}
1{
  temp_speed(ID,T2,S-1);
  temp_speed(ID,T2,S);
  temp_speed(ID,T2,S+1)
}1
:- chosen(ID,T,A,_,_,_,_,T2),
   A != wait,
   temp_speed(ID,T,S),
   max_speed(ID,MS),
   S-1 >= MS,
   S+1 <= 4.
\end{verbatim}

To account for the upper and lower bounds of a train's speed capabilities, three additional choice rules handle edge cases: when a train is at its maximum speed, when it is at its minimum speed, or when both conditions apply simultaneously. In the later case, trains are forced to maintain their current speed.

\begin{verbatim}
% Lower bound
1{
  temp_speed(ID,T2,S-1);
  temp_speed(ID,T2,S)
}1
:- chosen(ID,T,A,_,_,_,_,T2),
   A != wait,
   temp_speed(ID,T,S),
   max_speed(ID,MS),
   S-1 >= MS,
   S+1 > 4.

% Upper bound
1{
  temp_speed(ID,T2,S);
  temp_speed(ID,T2,S+1)
}1
:- chosen(ID,T,A,_,_,_,_,T2),
   A != wait,
   temp_speed(ID,T,S),
   max_speed(ID,MS),
   S-1 < MS,
   S+1 <= 4.

% In between
temp_speed(ID,T2,S)
:- chosen(ID,T,A,_,_,_,_,T2),
   A != wait,
   temp_speed(ID,T,S),
   max_speed(ID,MS),
   S-1 < MS,
   S+1 > 4.
\end{verbatim}

\paragraph{Post-Wait Acceleration}
When a train decides to wait, it loses its momentum. The encoding dictates that after taking a \texttt{wait} action, a train must resume movement at the lowest possible speed, which corresponds to an inverse speed of 4.

\begin{verbatim}
% after waiting, you have to start with speed 1/4
temp_speed(ID,T2,4)  :- chosen(ID,T,wait,_,_,_,_,T2).

% to be able to wait, you need to have speed 1/4
reachable(ID,T,wait,X,Y,T+1,Dir) :- position(ID,(X,Y),T,Dir), time(T+1), temp_speed(ID,T,4).
\end{verbatim}

\paragraph{Turning and Arrival Physics}
Realistic physics dictate that trains must decelerate before taking sharp turns or arriving at their destination. The model enforces that a train cannot take a turn at high speeds; it must slow down to an inverse speed of at least 3. Similarly, it must fully decelerate to an inverse speed of 4 upon reaching its final target.

\begin{verbatim}
% you have to take a turn with inverse speed >= 3 
:- position(ID,_,T,Dir1),
   temp_speed(ID,T,S),
   T2 = T + S,
   position(ID,_,T2,Dir2),
   Dir1 != Dir2,
   S < 3.

% you can only arrive at your goal if your inverse speed was 4 before
:- reached(ID,T2),
   temp_speed(ID,T1,S),
   T2 = T1 + S,
   S < 4.
\end{verbatim}

\paragraph{Updated Energy Costs}
The introduction of acceleration changes how energy consumption is calculated compared to the energy encoding, that does not take acceleration into account. The cost equations now heavily depend on the current inverse speed \texttt{S} and apply scaling multipliers to moving, waiting, and turning.

\begin{verbatim}
% Energy  cost of waiting
energy(ID,T+1,E) :- action(train(ID),wait,T),
                    not prev_wait(ID,T),
                    cost_fac(ID,CF),
                    E = -1 * 100 * 4/10 * CF * 1/4 * 1/4.

% Energy  cost of moving
energy(ID,T2,E) :- action(train(ID),A,T),
                   A != wait,
                   position(ID,_,T,Dir),
                   temp_speed(ID,T,S),
                   T2 = T + S,
                   position(ID,_,T2,Dir),
                   cost_fac(ID,CF),
                   E = 100 * CF * 1/S * 1/S.

% Energy  cost of turning
energy(ID,T2,E) :- position(ID,_,T,Dir1),
                   temp_speed(ID,T,S),
                   T2 = T + S,
                   position(ID,_,T2,Dir2),
                   Dir1 != Dir2,
                   cost_fac(ID,CF),
                   E = 100 * CF * 1/S * 1/S  * 11/10.
\end{verbatim}

\paragraph{Solver Heuristics}
To guide the solver's search process through the expanded state space, the encoding introduces specific \texttt{\#heuristic} directives to prioritize specific branches by penalizing waiting, scheduling trains one by one based on departure times, and instructing the solver to prefer paths where trains gain speed quickly.

\begin{verbatim}
% prefer to not wait
wait_position(ID,X2,Y2,T2,Dir2) 
:- reachable(ID,T,wait,X2,Y2,T2,Dir2).

#heuristic position(ID,(X2,Y2),T2,Dir2)
           : wait_position(ID,X2,Y2,T2,Dir2). [-50,1]

% schedule trains one by one
#heuristic position(ID,(X,Y),T,Dir)
           : start(ID,_,Dep,_). [-Dep,1]

% prefer getting faster
#heuristic temp_speed(ID,T,S)
           : position(ID,_,T,_). [-S,1]
\end{verbatim}

\subsection{Results}
For all trial runs a single environment was used. It consists of a 25 $\times$ 25 network with two stations and four agents. Three of the agents move at the speed of 1, the remaining one at speed 4.
\begin{figure}[htbp]
  \centering
  \includegraphics[width=0.8\textwidth]{environment.png}
  \caption{Flatland environment used for all experiments}
  \label{fig:environment_image}
\end{figure}

The baseline encoding was only used to establish rules for bringing trains from a starting point to an end point with a proper behavior for a given environment. That goal was reached successfully: The baseline encoding successfully generated paths for all four agents from their starting point to their end points.


\section{Conclusion}

\appendix

\section{Tables for Baseline Encoding}

\begin{table}[htbp]
  \centering
  \caption{Paths taken by agent 0 in the baseline encoding}
  \label{tab:baseline0}
  \begin{tabular}{llllll}
    \toprule
    agent & timestep & position & direction & status & given\_command \\
    \midrule
    0 & 0 & None & w & waiting & wait \\
    0 & 1 & None & w & waiting & wait \\
    0 & 2 & None & w & waiting & wait \\
    0 & 3 & None & w & waiting & wait \\
    0 & 4 & None & w & waiting & wait \\
    0 & 5 & None & w & waiting & wait \\
    0 & 6 & None & w & waiting & wait \\
    0 & 7 & None & w & waiting & wait \\
    0 & 8 & None & w & waiting & wait \\
    0 & 9 & None & w & waiting & wait \\
    0 & 10 & None & w & waiting & wait \\
    0 & 11 & None & w & waiting & wait \\
    0 & 12 & None & w & waiting & wait \\
    0 & 13 & None & w & waiting & wait \\
    0 & 14 & None & w & waiting & wait \\
    0 & 15 & None & w & waiting & wait \\
    0 & 16 & None & w & waiting & wait \\
    0 & 17 & None & w & waiting & wait \\
    0 & 18 & None & w & waiting & wait \\
    0 & 19 & None & w & waiting	 & wait \\
    0 & 20 & None & w & waiting & wait \\
    0 & 21 & None & w & waiting & wait \\
    0 & 22 & None & w & waiting & wait \\
    0 & 23 & None & w & waiting & wait \\
    0 & 24 & None & w & ready to depart & move\_forward \\
    0 & 25 & (7, 6) & w & moving & move\_forward \\
    0 & 26 & (7, 5) & w & moving & move\_forward \\
    0 & 27 & (8, 5) & s & moving & move\_forward \\
    0 & 28 & (8, 4) & w & moving & move\_forward \\
    0 & 29 & (8, 3) & w & moving & move\_forward \\
    0 & 30 & (8, 2) & w & moving & move\_forward \\
    0 & 31 & (8, 1) & w & moving & move\_forward \\
    0 & 32 & (9, 1) & s & moving & move\_forward \\
    0 & 33 & (10, 1) & s & moving & move\_forward \\
    0 & 34 & (11, 1) & s & moving & move\_forward \\
    0 & 35 & (12, 1) & s & moving & move\_forward \\
    0 & 36 & (13, 1) & s & moving & move\_forward \\
    0 & 37 & (13, 2) & e & moving & move\_forward \\
    0 & 38 & (13, 3) & e & moving & move\_forward \\
    0 & 39 & (13, 4) & e & moving & move\_forward \\
    0 & 40 & (13, 5) & e & moving & move\_forward \\
    0 & 41 & (13, 6) & e & moving & move\_forward \\
    0 & 42 & (13, 7) & e & moving & move\_forward \\
    0 & 43 & (13, 8) & e & moving & move\_forward \\
    0 & 44 & (13, 9) & e & moving & move\_forward \\
    0 & 45 & (13, 10) & e & moving & move\_forward \\
    0 & 46 & (13, 11) & e & moving & move\_forward \\
    0 & 47 & (13, 12) & e & moving & move\_forward \\
    0 & 48 & (12, 12) & n & moving & move\_right \\
    0 & 50 & (12, 14) & e & moving & move\_right \\
    0 & 51 & (13, 14) & s & moving & move\_forward \\
    0 & 52 & (13, 15) & e & moving & move\_forward \\
    0 & 53 & (13, 16) & e & moving & move\_forward \\
    \bottomrule
  \end{tabular}
\end{table}

\begin{table}[htbp]
  \centering
  \caption{Paths taken by agent 1 in the baseline encoding}
  \label{tab:baseline0}
  \begin{tabular}{llllll}
    \toprule
    agent & timestep & position & direction & status & given\_command \\
    \midrule
1 & 0 & None & e & waiting & wait \\
1 & 1 & None & e & waiting & wait \\
1 & 2 & None & e & waiting & wait \\
1 & 3 & None & e & waiting & wait \\
1 & 4 & None & e & waiting & wait \\
1 & 5 & None & e & waiting & wait \\
1 & 6 & None & e & waiting & wait \\
1 & 7 & None & e & waiting & wait \\
1 & 8 & None & e & waiting & wait \\
1 & 9 & None & e & waiting & wait \\
1 & 10 & None & e & waiting & wait \\
1 & 11 & None & e & waiting & wait \\
1 & 12 & None & e & waiting & wait \\
1 & 13 & None & e & waiting & wait \\
1 & 14 & None & e & waiting & wait \\
1 & 15 & None & e & waiting & wait \\
1 & 16 & None & e & waiting & wait \\
1 & 17 & None & e & waiting & wait \\
1 & 18 & None & e & waiting & wait \\
1 & 19 & None & e & waiting & wait \\
1 & 20 & None & e & waiting & wait \\
1 & 21 & None & e & ready to depart & move\_forward \\
1 & 22 & (12, 17) & e & moving & move\_forward \\
1 & 23 & (12, 18) & e & moving & move\_forward \\
1 & 24 & (12, 19) & e & moving & move\_forward \\
1 & 25 & (12, 20) & e & moving & move\_forward \\
1 & 26 & (12, 21) & e & moving & move\_forward \\
1 & 27 & (12, 22) & e & moving & move\_forward \\
1 & 28 & (11, 22) & n & moving & move\_forward \\
1 & 29 & (10, 22) & n & moving & move\_forward \\
1 & 30 & (9, 22)	 & n & moving & move\_forward \\
1 & 31 & (8, 22)	 & n & moving & move\_forward \\
1 & 32 & (7, 22)	 & n & moving & move\_forward \\
1 & 33 & (7, 21)	 & w & moving & move\_forward \\
1 & 34 & (7, 20)	 & w & moving & move\_forward \\
1 & 35 & (7, 19)	 & w & moving & move\_forward \\
1 & 36 & (7, 18)	 & w & moving & move\_forward \\
1 & 37 & (7, 17)	 & w & moving & move\_forward \\
1 & 38 & (7, 16)	 & w & moving & move\_forward \\
1 & 39 & (7, 15)	 & w & moving & move\_forward \\
1 & 40 & (7, 14)	 & w & moving & move\_forward \\
1 & 41 & (7, 13)	 & w & moving & move\_forward \\
1 & 42 & (7, 12)	 & w & moving & move\_left \\
1 & 43 & (8, 12)	 & s & moving & move\_forward \\
1 & 44 & (8, 11)	 & w & moving & move\_forward \\
1 & 45 & (8, 10)	 & w & moving & move\_forward \\
1 & 46 & (8, 9)	 & w & moving & move\_forward \\
1 & 47 & (8, 8)	 & w & moving & move\_forward \\
1 & 48 & (8, 7)	 & w & moving & move\_forward \\
    \bottomrule
  \end{tabular}
\end{table}

\begin{table}[htbp]
  \centering
  \caption{Paths taken by agent 2 in the baseline encoding}
  \label{tab:baseline0}
  \begin{tabular}{llllll}
    \toprule
    agent & timestep & position & direction & status & given\_command \\
    \midrule
2 & 0 & None & w & waiting & wait \\
2 & 1 & None & w & waiting & wait \\
2 & 2 & None & w & waiting & wait \\
2 & 3 & None & w & ready to depart & move\_forward \\
2 & 4 & (12, 17) & w & moving & move\_forward \\
2 & 8 & (12, 16) & w & moving & move\_forward \\
2 & 12 & (12, 15) & w & moving & move\_forward \\
2 & 16 & (12, 14) & w & moving & move\_forward \\
2 & 20 & (12, 13) & w & moving & move\_forward \\
2 & 24 & (12, 12) & w & moving & move\_forward \\
2 & 28 & (12, 11) & w & moving & move\_forward \\
2 & 32 & (11, 11) & n & moving & move\_forward \\
2 & 36 & (10, 11) & n & moving & wait \\
2 & 37 & (10, 11) & n & stopped & move\_forward \\
2 & 41 & (9, 11) & n & moving & move\_forward \\
2 & 45 & (8, 11) & n & moving & move\_left \\
2 & 49 & (8, 10) & w & moving & move\_forward \\
2 & 53 & (8, 9) & w & moving & move\_forward \\
2 & 54 & (8, 9) & w & moving & move\_forward \\
2 & 55 & (8, 9) & w & moving & move\_forward \\
    \bottomrule
  \end{tabular}
\end{table}

\begin{table}[htbp]
  \centering
  \caption{Paths taken by agent 3 in the baseline encoding}
  \label{tab:baseline0}
  \begin{tabular}{llllll}
    \toprule
    agent & timestep & position & direction & status & given\_command \\
    \midrule
3 & 0 & None & w & waiting & wait \\
3 & 1 & None & w & waiting & wait \\
3 & 2 & None & w & waiting & wait \\
3 & 3 & None & w & waiting & wait \\
3 & 4 & None & w & waiting & wait \\
3 & 5 & None & w & waiting & wait \\
3 & 6 & None & w & waiting & wait \\
3 & 7 & None & w & waiting & wait \\
3 & 8 & None & w & waiting & wait \\
3 & 9 & None & w & waiting & wait \\
3 & 10 & None & w & waiting & wait \\
3 & 11 & None & w & waiting & wait \\
3 & 12 & None & w & waiting & wait \\
3 & 13 & None & w & waiting & wait \\
3 & 14 & None & w & waiting & wait \\
3 & 15 & None & w & waiting & wait \\
3 & 16 & None & w & waiting & wait \\
3 & 17 & None & w & waiting & wait \\
3 & 18 & None & w & ready to depart & move\_forward \\
3 & 19 & (9, 6) & w & moving & move\_forward \\
3 & 20 & (9, 5) & w & moving & move\_forward \\
3 & 21 & (9, 4) & w & moving & move\_forward \\
3 & 22 & (9, 3) & w & moving & move\_forward \\
3 & 23 & (8, 3) & n & moving & move\_forward \\
3 & 24 & (8, 2) & w & moving & move\_forward \\
3 & 25 & (8, 1) & w & moving & move\_forward \\
3 & 26 & (9, 1) & s & moving & move\_forward \\
3 & 27 & (10, 1) & s & moving & move\_forward \\
3 & 28 & (11, 1) & s & moving & move\_forward \\
3 & 29 & (12, 1) & s & moving & move\_forward \\
3 & 30 & (13, 1) & s & moving & move\_forward \\
3 & 31 & (13, 2) & e & moving & move\_forward \\
3 & 32 & (13, 3) & e & moving & move\_forward \\
3 & 33 & (13, 4) & e & moving & move\_forward \\
3 & 34 & (13, 5) & e & moving & move\_forward \\
3 & 35 & (13, 6) & e & moving & move\_forward \\
3 & 36 & (13, 7) & e & moving & move\_forward \\
3 & 37 & (13, 8) & e & moving & move\_forward \\
3 & 38 & (13, 9) & e & moving & move\_forward \\
3 & 39 & (13, 10) & e & moving & move\_forward \\
3 & 40 & (13, 11) & e & moving & move\_forward \\
3 & 41 & (13, 12) & e & moving & move\_forward \\
3 & 42 & (12, 12) & n & moving & move\_right \\
3 & 43 & (12, 13) & e & moving & move\_forward \\
3 & 44 & (12, 14) & e & moving & move\_right \\
3 & 45 & (13, 14) & s & moving & move\_forward \\
3 & 46 & (13, 15) & e & moving & move\_forward \\
3 & 47 & (13, 16) & e & moving & move\_forward \\
    \bottomrule
  \end{tabular}
\end{table}

\section{Tables for Energy Encoding}

\end{document}